% !TeX root = ../../Book1.tex
\section{无界自伴算子的谱表示}
\label{b1:sec:2906}

有界自伴算子的谱积分在整个 Hilbert 空间上作用；无界情形首先要回答的是哪些向量允许代入。
本节从自伴算子的 Cayley 变换出发，把单位圆上的酉谱测度移到实轴，
并由谱测度同时恢复算子作用和定义域。随后把坐标函数换成一般 Borel 函数，
逐项核对闭性、伴随、乘积以及极限。特别地，两个算子的作用公式相同，
还不足以断言它们相等，必须同时核对定义域。

以下 \(H\) 为复 Hilbert 空间，内积对第一变量线性。
对投影值测度 \(E\) 及有界 Borel 函数 \(h\)，暂记
\[
 h(E)=\int h\,\mathrm dE.
\]
这是\secref{b1:sec:2904}已经构造的有界谱积分；这个记号不预先要求有一个无界算子。
标量正测度记为
\[
 \mu_x(B)=\langle E(B)x,x\rangle=\|E(B)x\|^2.
\]
本节的 Cayley 证明路线参见\cite[Theorem 13.30, pp.~368--369]{Rudin1991Functional}；
一般谱积分的定义域、乘积与伴随参见\cite[Theorem 13.24, pp.~362--364]{Rudin1991Functional}。

\subsection{Cayley 变换}
\label{b1:subsec:290601}

\cref{b1:thm:29-self-adjoint-criterion}表明，对稠密定义的自伴算子 \(A\)，
\(A+iI\) 和 \(A-iI\) 都把 \(D(A)\) 双射到 \(H\)。
因此可以先用 \((A+iI)^{-1}\) 把任意向量送入定义域，再作用 \(A-iI\)。

\begin{proposition}[自伴算子的 Cayley 变换]
\label{b1:prop:29-cayley}
设 \(H\) 为复 Hilbert 空间，\(A:D(A)\subseteq H\to H\) 为稠密定义的自伴算子。则
\[
 U=(A-iI)(A+iI)^{-1}
   =I-2i(A+iI)^{-1}
\]
是 \(H\) 上的酉算子，称为 \(A\) 的\term[cayley-bianhuan]{Cayley 变换}[Cayley transform]。
它满足
\begin{equation}\label{b1:eq:29-cayley-domain}
 I-U=2i(A+iI)^{-1},\qquad
 \ker(I-U)=\{0\},\qquad
 \operatorname{ran}(I-U)=D(A),
\end{equation}
并且
\begin{equation}\label{b1:eq:29-cayley-inverse}
 A(I-U)y=i(I+U)y\qquad(y\in H).
\end{equation}
\end{proposition}
\begin{proof}
由对称性，\(\langle Ax,x\rangle\) 为实数，故对 \(x\in D(A)\) 有
\[
 \|(A\pm iI)x\|^2=\|Ax\|^2+\|x\|^2.
\]
两算子满射，再由这个等式可知它们的逆作为 \(H\to H\) 的映射有界，范数至多为 \(1\)。
它们的值域都恰为 \(D(A)\)。若 \(y=(A+iI)x\)，则
\[
 Uy=(A-iI)x=y-2ix,
\]
所以 \(U=I-2i(A+iI)^{-1}\) 在整个 \(H\) 上有界，且
\(\|Uy\|=\|(A-iI)x\|=\|(A+iI)x\|=\|y\|\)。
因为 \(A-iI\) 满射，\(U\) 也满射，故为酉算子。

由 \(I-U=2i(A+iI)^{-1}\)，立即得到核与值域的两个结论。
仍把 \(y\) 写成 \((A+iI)x\)，则
\[
 (I-U)y=2ix,\qquad (I+U)y=2Ax.
\]
前式保证 \((I-U)y\in D(A)\)，并给出
\(A(I-U)y=2iAx=i(I+U)y\)。这里每一次作用 \(A\) 的向量都已确定属于定义域。
\end{proof}

记单位圆为 \(\mathbb T=\{z\in\mathbb C:|z|=1\}\)。两个互逆的连续映射
\begin{equation}\label{b1:eq:29-cayley-maps}
 \phi(t)=\frac{t-i}{t+i}\quad(t\in\mathbb R),\qquad
 \theta(z)=i\frac{1+z}{1-z}\quad(z\in\mathbb T\setminus\{1\})
\end{equation}
把实轴与去掉 \(1\) 的单位圆对应起来。
对 \(|z|=1\)、\(z\ne1\)，利用 \(\overline z=z^{-1}\) 可核对
\(\overline{\theta(z)}=\theta(z)\)，故 \(\theta\) 确实取实值。
还将用到
\begin{equation}\label{b1:eq:29-cayley-weight}
 \bigl(1+|\theta(z)|^2\bigr)\cdot|1-z|^2
 =|1-z|^2+|1+z|^2=4.
\end{equation}

由酉谱定理\cref{b1:thm:29-unitary-spectral}，存在唯一的投影值测度 \(F\)，使
\[
 U=\int_{\mathbb T}z\,\mathrm dF(z).
\]
有界谱积分的乘法公式给出
\((I-U)F(\{1\})=0\)。结合\eqref{b1:eq:29-cayley-domain}的单射性，得到
\begin{equation}\label{b1:eq:29-cayley-no-atom}
 F(\{1\})=0.
\end{equation}
这个结论只排除了 \(1\) 处的谱原子。若 \(D(A)\ne H\)，则 \(I-U\) 的值域是稠密真子空间，
因而 \(1\) 仍属于 \(U\) 的谱。逆变换 \(\theta\) 在该点附近无界，
这正是需要引入无界谱积分的地方。

\subsection{无界自伴谱定理}
\label{b1:subsec:290602}

先把圆上的投影值测度推到实轴：对 Borel 集 \(B\subseteq\mathbb R\)，定义
\begin{equation}\label{b1:eq:29-real-spectral-measure}
 E(B)=F\bigl(\theta^{-1}(B)\bigr).
\end{equation}
原像保持交与不交并，故 \(E\) 保持投影乘法和强可数可加性。
由\eqref{b1:eq:29-cayley-no-atom}，又有
\(E(\mathbb R)=F(\mathbb T\setminus\{1\})=I\)，所以 \(E\) 是投影值测度。
若 \(h\) 为实轴上的有界 Borel 函数，在 \(z=1\) 处把 \(h\circ\theta\) 任意补成有限值，则
\begin{equation}\label{b1:eq:29-cayley-pushforward}
 h(E)=\int_{\mathbb T}h(\theta(z))\,\mathrm dF(z).
\end{equation}
对示性函数，这就是定义；对简单函数，利用线性性；对一般有界 Borel 函数，
用一致有界的简单函数逐点逼近，再用有界谱积分的强收敛结论即可。

对于尚未定义的坐标积分，先取有界截断
\[
 \tau_n(t)=t\cdot\mathbf1_{[-n,n]}(t),\qquad A_n=\tau_n(E).
\]
若 \(\int t^2\,\mathrm d\mu_x(t)<\infty\)，则对 \(m\ge n\)，
\[
 \|(A_m-A_n)x\|^2
 =\int_{\{n<|t|\le m\}}t^2\,\mathrm d\mu_x(t)\longrightarrow0.
\]
因此 \(A_nx\) 在 \(H\) 中有极限。下面证明，这个平方可积条件恰好刻画原来的定义域。

\begin{theorem}[无界自伴谱定理]
\label{b1:thm:29-unbounded-spectral}
设 \(H\) 为复 Hilbert 空间，\(A:D(A)\subseteq H\to H\) 为稠密定义的自伴算子。
则在 \(\mathbb R\) 的 Borel 集上存在唯一的投影值测度 \(E\)，使
\begin{equation}\label{b1:eq:29-spectral-domain}
 D(A)=\left\{x\in H:\int_{\mathbb R}t^2\,\mathrm d\mu_x(t)<\infty\right\},
 \qquad \mu_x(B)=\langle E(B)x,x\rangle,
\end{equation}
并且
\begin{equation}\label{b1:eq:29-spectral-action}
 Ax=\lim_{n\to\infty}\int_{[-n,n]}t\,\mathrm dE(t)x,
 \qquad
 \|Ax\|^2=\int_{\mathbb R}t^2\,\mathrm d\mu_x(t)
 \quad(x\in D(A)).
\end{equation}
对每个 \(w\in\mathbb C\setminus\mathbb R\)，\(A-wI:D(A)\to H\) 为双射，其逆满足
\begin{equation}\label{b1:eq:29-unbounded-resolvent}
 (A-wI)^{-1}=\int_{\mathbb R}\frac1{t-w}\,\mathrm dE(t),\qquad
 \|(A-wI)^{-1}\|\le\frac1{|\operatorname{Im}w|}.
\end{equation}
\end{theorem}
\begin{proof}
若 \(H=\{0\}\)，则零投影值测度直接给出全部结论。以下设 \(H\ne\{0\}\)。
取由\eqref{b1:eq:29-real-spectral-measure}构造的 \(E\)。
先证明 \(D(A)\) 中的向量满足平方可积条件。
由\eqref{b1:eq:29-cayley-domain}，可把任意 \(x\in D(A)\) 唯一写成
\(x=(I-U)y\)。记 \(\nu_y(C)=\langle F(C)y,y\rangle\)。
对 Borel 集 \(C\subseteq\mathbb T\)，有界谱积分的乘法与范数公式给出
\[
 \nu_x(C)=\|F(C)(I-U)y\|^2
 =\int_C|1-z|^2\,\mathrm d\nu_y(z).
\]
推前标量测度，再用\eqref{b1:eq:29-cayley-weight}，得到
\begin{equation}\label{b1:eq:29-cayley-domain-norm}
 \int_{\mathbb R}(1+t^2)\,\mathrm d\mu_x(t)
 =\int_{\mathbb T\setminus\{1\}}
   (1+|\theta(z)|^2)\cdot|1-z|^2\,\mathrm d\nu_y(z)
 =4\cdot\|y\|^2.
\end{equation}
这里无界非负函数的换元可先对简单函数验证，再由单调收敛定理得到。
特别地，\(x\) 满足所需的平方可积条件。另一方面，
\[
 A_nx
 =i\int_{\{z\ne1:\,|\theta(z)|\le n\}}
       (1+z)\,\mathrm dF(z)y
 \longrightarrow i(I+U)y=Ax.
\]
收敛来自有界谱积分的支配收敛，因为乘数的绝对值至多为 \(2\)，
而这些集合递增覆盖 \(\mathbb T\setminus\{1\}\)。

反过来，设 \(x\in H\) 满足 \(\int t^2\,\mathrm d\mu_x<\infty\)，令
\[
 y_n=\left(\frac{t+i}{2i}\cdot\mathbf1_{[-n,n]}(t)\right)(E)x.
\]
由于 \(|(t+i)/(2i)|^2=(1+t^2)/4\) 对 \(\mu_x\) 可积，
同样的范数差估计表明 \((y_n)\) 为 Cauchy 列。设其极限为 \(y\)。
由\eqref{b1:eq:29-cayley-pushforward}和 \(F(\{1\})=0\)，有 \(U=\phi(E)\)，故
\[
 (I-U)y_n
 =\left(\left(1-\frac{t-i}{t+i}\right)
          \frac{t+i}{2i}\cdot\mathbf1_{[-n,n]}(t)\right)(E)x
 =E([-n,n])x\longrightarrow x.
\]
\(I-U\) 有界，因而 \((I-U)y=x\)。
由其值域等于 \(D(A)\)，得到 \(x\in D(A)\)。这证明了定义域等式和作用公式。
最后，\(A_nx\to Ax\) 及单调收敛给出
\[
 \|Ax\|^2=\lim_n\|A_nx\|^2
 =\lim_n\int_{[-n,n]}t^2\,\mathrm d\mu_x
 =\int_{\mathbb R}t^2\,\mathrm d\mu_x.
\]

现在取 \(w\notin\mathbb R\)，置 \(r_w(t)=(t-w)^{-1}\)、\(R=r_w(E)\)。
有
\[
 |r_w(t)|\le\frac1{|\operatorname{Im}w|},\qquad
 |t r_w(t)|\le1+\frac{|w|}{|\operatorname{Im}w|}.
\]
对 \(v\in H\)，有界谱积分给出的标量测度关系为
\[
 \mu_{Rv}(B)=\int_B|r_w(t)|^2\,\mathrm d\mu_v(t).
\]
因此 \(\int t^2\,\mathrm d\mu_{Rv}<\infty\)，即 \(Rv\in D(A)\)。
再由坐标截断和有界谱积分的强收敛，
\[
 ARv=\lim_n(\tau_n r_w)(E)v=(t r_w)(E)v,
 \qquad (A-wI)Rv=v.
\]
对 \(x\in D(A)\)，\(\langle Ax,x\rangle\) 为实数，故
\[
 |\operatorname{Im}w|\cdot\|x\|^2
 =\bigl|\operatorname{Im}\langle(A-wI)x,x\rangle\bigr|
 \le\|(A-wI)x\|\cdot\|x\|.
\]
所以 \(A-wI\) 单射。结合已经证明的右逆，它为双射，且逆的范数满足所述估计。

最后证明谱测度的唯一性。设另一个实轴投影值测度 \(E'\) 也满足
\eqref{b1:eq:29-spectral-domain}和\eqref{b1:eq:29-spectral-action}。
上面的预解算子论证只用了这两个性质以及有界谱积分规则，所以同样适用于 \(E'\)。
在 \(w=-i\) 处得到
\[
 (A+iI)^{-1}=\left(\frac1{t+i}\right)(E'),\qquad U=\phi(E').
\]
在单位圆上定义 \(F'(C)=E'(\phi^{-1}(C))\)。
它是投影值测度，满足 \(F'(\{1\})=0\)，而且
\(U=\int_{\mathbb T}z\,\mathrm dF'(z)\)。
酉谱定理的唯一性给出 \(F'=F\)。
对任意 Borel 集 \(B\subseteq\mathbb R\)，\(\theta^{-1}(B)=\phi(B)\)
是 \(\mathbb T\setminus\{1\}\) 的 Borel 集，故
\[
 E'(B)=F'(\theta^{-1}(B))=F(\theta^{-1}(B))=E(B).
\]
\end{proof}

这个定理中的唯一性同时包含定义域；只在一个较小的稠密子空间上给出积分公式，
不能代替\eqref{b1:eq:29-spectral-domain}的等式。
与有界情形一样，记
\[
 \rho(A)=\{w\in\mathbb C:A-wI:D(A)\to H
       \text{ 为双射且逆在 }H\text{ 上有界}\},\qquad
 \sigma(A)=\mathbb C\setminus\rho(A).
\]
由于 \(A\) 闭，双射时逆的有界性也可由闭图像定理推出。
\eqref{b1:eq:29-unbounded-resolvent}说明 \(\sigma(A)\subseteq\mathbb R\)。
谱测度还精确说明哪些实数属于谱。

\begin{proposition}[谱测度的支集与谱]
\label{b1:prop:29-unbounded-spectral-support}
设 \(H\ne\{0\}\) 为复 Hilbert 空间，\(A:D(A)\subseteq H\to H\) 为稠密定义的自伴算子，
\(E\) 为\cref{b1:thm:29-unbounded-spectral}给出的实轴谱测度。则
\[
 \sigma(A)=\{s\in\mathbb R:
 E((s-\varepsilon,s+\varepsilon))\ne0\text{ 对每个 }\varepsilon>0\text{ 成立}\}.
\]
这个集合为非空闭集，且 \(E(\sigma(A))=I\)。
\end{proposition}
\begin{proof}
把右端记为 \(K\)。它的补集为开集。
若 \(s\notin K\)，取 \(\varepsilon>0\) 使 \(E((s-\varepsilon,s+\varepsilon))=0\)，并定义
\[
 r(t)=
 \begin{cases}
  (t-s)^{-1},&|t-s|\ge\varepsilon,\\
  0,&|t-s|<\varepsilon.
 \end{cases}
\]
\(r\) 与 \(tr\) 都有界。与预解公式的证明相同，标量测度范数公式先给出
\(r(E)H\subseteq D(A)\)，坐标截断再给出
\[
 (A-sI)r(E)v=((t-s)r)(E)v=v\qquad(v\in H).
\]
对 \(x\in D(A)\)，还可在另一侧计算
\[
 \begin{aligned}
 r(E)(A-sI)x
 &=\lim_n r(E)(A_n-sI)x\\
 &=\bigl(tr-sr\bigr)(E)x=x.
 \end{aligned}
\]
第二行的极限成立，是因为 \(r\tau_n\) 一致有界并逐点趋于 \(rt\)，
而 \((t-s)r=1\) 在 \(E\)-几乎处处成立。因此 \(s\in\rho(A)\)。

若 \(s\in K\)，则对每个 \(0<\varepsilon<1\)，可在
\(E((s-\varepsilon,s+\varepsilon))H\) 中选一个单位向量 \(x_\varepsilon\)。
它的标量谱测度集中在这个有界区间上，故 \(x_\varepsilon\in D(A)\)，且
\[
 \|(A-sI)x_\varepsilon\|^2
 =\int_{(s-\varepsilon,s+\varepsilon)}|t-s|^2
          \,\mathrm d\mu_{x_\varepsilon}(t)
 \le\varepsilon^2.
\]
这里在有界区间上，\(A\) 的作用就是有界坐标积分。
若 \((A-sI)^{-1}\) 存在且有界，就会有
\(1\le\|(A-sI)^{-1}\|\cdot\varepsilon\)，与 \(\varepsilon\downarrow0\) 矛盾。
所以 \(s\in\sigma(A)\)，得到 \(K=\sigma(A)\)。

对每个 \(s\notin K\)，可在其一个零投影邻域内选取包含 \(s\) 的有理端点开区间。
所有这些区间构成可数族，覆盖 \(\mathbb R\setminus K\)，且各自的 \(E\)-投影为零。
对每个 \(x\in H\)，标量测度的可数次可加性给出
\(\mu_x(\mathbb R\setminus K)=0\)，所以 \(E(\mathbb R\setminus K)=0\)、\(E(K)=I\)。
若 \(K\) 为空，这会与 \(E(\mathbb R)=I\ne0\) 矛盾。
\end{proof}

\subsection{无界函数演算的定义域}
\label{b1:subsec:290603}

现在固定实轴上的投影值测度 \(E\)。下面的构造只依赖 \(E\)，
因而也可以用于尚未指定自伴算子的任意这种投影值测度。
所谓 \(E\)-几乎处处，是指除去一个满足 \(E(N)=0\) 的 Borel 集 \(N\)；
这样的集合对每个 \(\mu_x\) 都为零测集。
本节只考虑处处取有限复数值的 Borel 函数，允许函数无界。

\begin{definition}[无界谱积分与最大定义域]
\label{b1:def:29-unbounded-calculus}
设 \(H\) 为复 Hilbert 空间，\(E\) 为实轴 Borel 集上的投影值测度，
\(\mu_x(B)=\langle E(B)x,x\rangle\)。设 \(f:\mathbb R\to\mathbb C\) 为有限值 Borel 函数，记
\[
 C_n=\{t:|f(t)|\le n\},\qquad P_n=E(C_n),\qquad
 f_n=f\cdot\mathbf1_{C_n}.
\]
定义谱积分 \(T_f=\int f\,\mathrm dE\) 的\term[zuida-dingyiyu]{最大定义域}[maximal domain]为
\begin{equation}\label{b1:eq:29-borel-maximal-domain}
 D_f=D(T_f)=\left\{x\in H:\int_{\mathbb R}|f(t)|^2\,\mathrm d\mu_x(t)<\infty\right\},
\end{equation}
并对 \(x\in D_f\) 置
\begin{equation}\label{b1:eq:29-borel-truncation}
 T_fx=\lim_{n\to\infty}f_n(E)x.
\end{equation}
若 \(E\) 为自伴算子 \(A\) 的谱测度，则记
\[
 f(A)=T_f=\int_{\mathbb R}f(t)\,\mathrm dE(t).
\]
\symidx{\(D(f(A))\)：无界谱函数演算的最大定义域}
\end{definition}

先验证定义确实给出稠密定义的线性算子。
由有界谱积分的范数公式及单调收敛，
\begin{equation}\label{b1:eq:29-borel-domain-test}
 \sup_n\|f_n(E)x\|^2=\int|f|^2\,\mathrm d\mu_x.
\end{equation}
所以 \(D_f\) 恰为使左侧有限的向量全体，算子范数的三角不等式说明它为线性子空间。
对 \(x\in D_f\)、\(m\ge n\)，有
\[
 \|(f_m(E)-f_n(E))x\|^2
 =\int_{\{n<|f|\le m\}}|f|^2\,\mathrm d\mu_x\longrightarrow0,
\]
故\eqref{b1:eq:29-borel-truncation}中的极限存在。
各个截断算子的线性性再说明 \(T_f\) 线性，且
\begin{equation}\label{b1:eq:29-borel-norm}
 \|T_fx\|^2=\int|f|^2\,\mathrm d\mu_x\qquad(x\in D_f).
\end{equation}
对任意 \(y\in H\)，\(P_ny\) 的标量谱测度集中在 \(C_n\) 上，因而
\(\int|f|^2\,\mathrm d\mu_{P_ny}\le n^2\cdot\|y\|^2\)，即 \(P_nH\subseteq D_f\)。
\(f\) 处处有限意味着 \(C_n\uparrow\mathbb R\)，从而 \(P_n\to I\) 强收敛，
\(D_f\) 在 \(H\) 中稠密。

若 \(f\) 有界，则充分大的 \(n\) 满足 \(f_n=f\)，所以此定义与原来的有界谱积分一致。
若 \(f\) 和 \(g\) 在 \(E\)-几乎处处相等，则各自的平方可积域相同，
截断函数的有界谱积分也相同，故 \(T_f=T_g\)。
对于坐标函数 \(f(t)=t\)，\cref{b1:thm:29-unbounded-spectral}说明得到的算子恰为 \(A\)，
包括定义域的等式。

\begin{proposition}[谱投影与标量谱测度的变换]
\label{b1:prop:29-unbounded-spectral-commutation}
设 \(H\) 为复 Hilbert 空间，\(E\) 为实轴 Borel 集上的投影值测度，
\(f\) 为有限值 Borel 函数，\(T_f=\int f\,\mathrm dE\) 取最大定义域 \(D_f\)。
设 \(B\subseteq\mathbb R\) 为 Borel 集。
则对 \(x\in D_f\)，有
\[
 E(B)x\in D_f,\qquad T_fE(B)x=E(B)T_fx,
\]
且
\begin{equation}\label{b1:eq:29-unbounded-scalar-transport}
 \mu_{T_fx}(B)=\int_B|f|^2\,\mathrm d\mu_x.
\end{equation}
若 \(f\) 在 \(B\) 上有界，则 \(E(B)H\subseteq D_f\)，并且作为整个 \(H\) 上的有界算子，
\[
 T_fE(B)=(f\cdot\mathbf1_B)(E).
\]
\end{proposition}
\begin{proof}
投影乘法给出
\(\mu_{E(B)x}(C)=\mu_x(B\cap C)\)。因此
\[
 \int|f|^2\,\mathrm d\mu_{E(B)x}=\int_B|f|^2\,\mathrm d\mu_x<\infty.
\]
每个有界截断 \(f_n(E)\) 都与 \(E(B)\) 交换，令 \(n\to\infty\) 即得作用公式。
再用\eqref{b1:eq:29-borel-norm}，
\[
 \mu_{T_fx}(B)=\|E(B)T_fx\|^2
 =\|T_fE(B)x\|^2=\int_B|f|^2\,\mathrm d\mu_x.
\]
若 \(|f|\le M\) 于 \(B\)，上述积分对任意 \(x\in H\) 至多为 \(M^2\cdot\|x\|^2\)。
并且对 \(n\ge M\)，\(f_n\cdot\mathbf1_B=f\cdot\mathbf1_B\)，所以截断定义直接给出最后的等式。
\end{proof}

\begin{theorem}[无界谱积分的闭性与伴随]
\label{b1:thm:29-unbounded-calculus-adjoint}
设 \(H\) 为复 Hilbert 空间，\(E\) 为实轴 Borel 集上的投影值测度，
\(f\) 为有限值 Borel 函数，\(T_f=\int f\,\mathrm dE\) 取最大定义域 \(D_f\)。
则 \(T_f\) 是稠密定义的闭算子，且
\begin{equation}\label{b1:eq:29-unbounded-calculus-adjoint}
 T_f^*=T_{\overline f},\qquad
 D(T_f^*)=D(T_{\overline f})=D_f.
\end{equation}
特别地，若 \(f\) 在 \(E\)-几乎处处为实值，则 \(T_f\) 自伴。
\end{theorem}
\begin{proof}
稠密定义已经验证。先直接证明闭性。
设 \(x_k\in D_f\)、\(x_k\to x\)、\(T_fx_k\to y\)。
固定 \(n\)，由上一命题有
\[
 f_n(E)x_k=P_nT_fx_k.
\]
两侧有界算子均可通过极限，得到 \(f_n(E)x=P_ny\)。于是
\(\|f_n(E)x\|\le\|y\|\) 对所有 \(n\) 成立。
\eqref{b1:eq:29-borel-domain-test}说明 \(x\in D_f\)，再令 \(n\to\infty\)，得到
\(T_fx=\lim_nP_ny=y\)，证明了闭性。

由于 \(|\overline f|=|f|\)，有 \(D_{\overline f}=D_f\)。
对 \(x\in D_f\)、\(y\in D_{\overline f}\)，有界谱积分的伴随公式给出
\[
 \langle T_fx,y\rangle
 =\lim_n\langle f_n(E)x,y\rangle
 =\lim_n\langle x,\overline{f_n}(E)y\rangle
 =\langle x,T_{\overline f}y\rangle.
\]
因此 \(T_{\overline f}\subseteq T_f^*\)。

反过来，设 \(y\in D(T_f^*)\)、\(v=T_f^*y\)。
对任意 \(u\in H\)，\(P_nu\in D_f\)，所以
\[
 \langle f_n(E)u,y\rangle
 =\langle T_fP_nu,y\rangle
 =\langle P_nu,v\rangle
 =\langle u,P_nv\rangle.
\]
由有界谱积分的伴随公式及表示向量唯一性，
\(\overline{f_n}(E)y=P_nv\)。因此
\[
 \int_{C_n}|f|^2\,\mathrm d\mu_y
 =\|\overline{f_n}(E)y\|^2\le\|v\|^2.
\]
单调收敛给出 \(y\in D_{\overline f}\)。令 \(n\to\infty\)，再得到
\(T_{\overline f}y=v\)，从而伴随等式连同定义域等式全部成立。
最后，若 \(f=\overline f\) 于 \(E\)-几乎处处，两个谱积分相同，故 \(T_f=T_f^*\)。
\end{proof}

上述闭性说明最大定义域对图像极限封闭。
若再把 \(T_f\) 限制到一个较小的稠密子空间，所得算子可能不闭，也可能不再自伴。
因此，谱积分中的平方可积条件是算子本身的一部分。

\subsection{谱演算的乘积与极限}
\label{b1:subsec:290604}

以下把记号恢复为 \(f(A)\)。所有函数都相对于同一个自伴算子 \(A\) 的谱测度计算。
算子乘积先作用右侧因子，约定
\[
 D(f(A)g(A))
 =\{x\in D(g(A)):g(A)x\in D(f(A))\}.
\]
算子和的定义域则是 \(D(f(A))\cap D(g(A))\)。

\begin{proposition}[乘积和加法的定义域]
\label{b1:prop:29-unbounded-calculus-products}
设 \(H\) 为复 Hilbert 空间，\(A:D(A)\subseteq H\to H\) 为稠密定义的自伴算子，
\(f,g:\mathbb R\to\mathbb C\) 为有限值 Borel 函数，函数演算均取最大定义域。则
\begin{equation}\label{b1:eq:29-product-domain}
 D(f(A)g(A))=D(g(A))\cap D((fg)(A)),
\end{equation}
并且
\begin{equation}\label{b1:eq:29-product-restriction}
 f(A)g(A)x=(fg)(A)x
 \qquad\bigl(x\in D(f(A)g(A))\bigr).
\end{equation}
对于加法，有
\begin{equation}\label{b1:eq:29-sum-restriction}
 D(f(A))\cap D(g(A))\subseteq D((f+g)(A)),\qquad
 (f(A)+g(A))x=(f+g)(A)x
\end{equation}
在左端定义域上成立。上述两个算子限制均可闭，且
\begin{equation}\label{b1:eq:29-product-sum-closures}
 \overline{f(A)g(A)}=(fg)(A),\qquad
 \overline{f(A)+g(A)}=(f+g)(A).
\end{equation}
\end{proposition}
\begin{proof}
对 \(x\in D(g(A))\)，由\eqref{b1:eq:29-unbounded-scalar-transport}，先对简单非负函数，
再用单调收敛，可得
\[
 \int|f|^2\,\mathrm d\mu_{g(A)x}
 =\int|fg|^2\,\mathrm d\mu_x.
\]
所以 \(g(A)x\in D(f(A))\) 当且仅当 \(x\in D((fg)(A))\)，证明了精确的乘积域。

令
\[
 B_n=\{|f|\le n,\ |g|\le n\},\qquad Q_n=E(B_n).
\]
两函数都取有限值，故 \(B_n\uparrow\mathbb R\)、\(Q_n\to I\) 强收敛。
在 \(B_n\) 上两个函数有界，可用有界谱积分的乘法公式。
对乘积域中的 \(x\)，上一小节的谱投影交换关系给出
\[
 \begin{aligned}
 Q_nf(A)g(A)x
 &=f(A)Q_ng(A)x\\
 &=(f\cdot\mathbf1_{B_n})(E)(g\cdot\mathbf1_{B_n})(E)x\\
 &=(fg\cdot\mathbf1_{B_n})(E)x
 =Q_n(fg)(A)x.
 \end{aligned}
\]
令 \(n\to\infty\) 得\eqref{b1:eq:29-product-restriction}。

若 \(x\in D(f(A))\cap D(g(A))\)，由
\(|f+g|^2\le2\cdot|f|^2+2\cdot|g|^2\)，首先有 \(x\in D((f+g)(A))\)。
再在同一 \(B_n\) 上使用有界谱积分的加法公式，得到
\[
 Q_n(f(A)x+g(A)x)
 =((f+g)\cdot\mathbf1_{B_n})(E)x
 =Q_n(f+g)(A)x.
\]
令 \(n\to\infty\)，得到加法的作用公式。

最后，这两个限制都包含 \(\bigcup_nQ_nH\) 于定义域中，故稠密定义；
又分别包含于闭算子 \((fg)(A)\) 与 \((f+g)(A)\)，因而可闭。
分别取 \(h=fg\) 或 \(h=f+g\)。对每个 \(x\in D(h(A))\)，置 \(x_n=Q_nx\)。
则 \(x_n\) 属于相应限制的定义域，并且
\[
 \|x-x_n\|^2+\|h(A)x-h(A)x_n\|^2
 =\int_{B_n^c}(1+|h|^2)\,\mathrm d\mu_x\longrightarrow0.
\]
这里被积函数对 \(\mu_x\) 可积，且 \(B_n\uparrow\mathbb R\)。
所以每个 \((x,h(A)x)\) 都在相应限制的图像闭包中。
反向包含由 \(h(A)\) 的闭性保证，证明了两个闭包等式。
\end{proof}

由\eqref{b1:eq:29-product-domain}，当右侧因子 \(g\) 有界时，确有
\(f(A)g(A)=(fg)(A)\)。仅知道左侧因子 \(f\) 有界，则一般只能得到扩张关系
\(f(A)g(A)\subseteq(fg)(A)\)。加法中的域包含也可能为严格包含。

在 \(H=\ell^2(\mathbb N)\) 上，可以把这两个现象直接写出。
定义投影值测度
\[
 E(B)(x_1,x_2,\ldots)
 =(\mathbf1_B(1)\cdot x_1,\mathbf1_B(2)\cdot x_2,\ldots).
\]
其强可数可加性由每个向量的平方可和性及标量支配收敛得到，且
\(\mu_x=\sum_{k\ge1}|x_k|^2\cdot\delta_k\)。
对坐标函数应用上一小节的谱积分，得到自伴算子
\[
 Mx=(kx_k)_{k\ge1},\qquad
 D(M)=\left\{x\in\ell^2:\sum_{k\ge1}k^2\cdot|x_k|^2<\infty\right\}.
\]
由无界自伴谱定理的唯一性，它的谱测度正是这个 \(E\)。
若 \(g(t)=t\)、\(f(t)=(1+|t|)^{-1}\)，则 \(f(M)\) 与 \((fg)(M)\) 都在整个空间上有界，
而 \(f(M)g(M)\) 的定义域仍为 \(D(M)\)。
例如 \(x=(1/k)_{k\ge1}\) 属于 \(\ell^2\) 却不属于 \(D(M)\)，所以这两个算子不相等。
若改取 \(f(t)=t\)、\(g(t)=-t\)，则 \(f(M)+g(M)\) 是定义在 \(D(M)\) 上的零算子，
而 \((f+g)(M)\) 是定义在整个 \(\ell^2\) 上的零算子。

\ProseParagraphBreak

函数列的点态极限也要连同可积的控制条件使用。

\begin{proposition}[谱积分的支配收敛]
\label{b1:prop:29-unbounded-calculus-dominated}
设 \(H\) 为复 Hilbert 空间，\(A:D(A)\subseteq H\to H\) 为稠密定义的自伴算子，
\(E\) 为其谱测度。设 \(f_n,f\) 为有限值 Borel 函数，\(h\ge0\) 为有限值 Borel 函数。
若在一个共同的 \(E\)-零测集之外，
\[
 f_n(t)\longrightarrow f(t),\qquad |f_n(t)|\le h(t)\quad(n\ge1),
\]
则对每个 \(x\in D(h(A))\)，均有 \(x\in D(f_n(A))\cap D(f(A))\)，且
\[
 f_n(A)x\longrightarrow f(A)x\quad\text{于 }H.
\]
\end{proposition}
\begin{proof}
取极限给出 \(|f|\le h\) 于共同零测集之外，故
\[
 \int|f_n|^2\,\mathrm d\mu_x\le\int h^2\,\mathrm d\mu_x<\infty,
 \qquad
 \int|f|^2\,\mathrm d\mu_x\le\int h^2\,\mathrm d\mu_x<\infty.
\]
由加法公式及范数公式，
\[
 \|f_n(A)x-f(A)x\|^2
 =\int|f_n-f|^2\,\mathrm d\mu_x.
\]
被积函数逐点趋于零，且由 \(4h^2\) 控制。
标量支配收敛定理因此给出右端趋于零。
\end{proof}

这个结论是对固定向量的范数收敛。
即使所有函数都一致有界，也不能据此推出算子范数收敛。
例如，对上面的 \(M\) 取 \(f_n=\mathbf1_{(n,\infty)}\)，有 \(f_n(M)x\to0\) 对每个 \(x\in\ell^2\) 成立，
但 \(f_n(M)e_{n+1}=e_{n+1}\)，所以 \(\|f_n(M)\|=1\) 对所有 \(n\) 成立。

最后把谱截断表述为定义域中的逼近。
这里使用图像范数下稠密的定义子空间（英文称 core）：对闭算子 \(T\)，
它指在 \(D(T)\) 中关于图像范数 \(\|x\|_T=\|x\|+\|Tx\|\) 稠密的线性子空间
\(\mathcal C\subseteq D(T)\)。
这等价于 \(T|_{\mathcal C}\) 的图像闭包等于 \(T\) 的图像。

\begin{proposition}[谱截断的图像范数逼近]
\label{b1:prop:29-spectral-graph-core}
设 \(H\) 为复 Hilbert 空间，\(A:D(A)\subseteq H\to H\) 为稠密定义的自伴算子，
\(E\) 为其谱测度，\(f\) 为有限值 Borel 函数，令 \(P_n=E(\{|f|\le n\})\)。
则 \(\mathcal C_f=\bigcup_{n\ge1}P_nH\) 在 \(D(f(A))\) 中关于图像范数稠密。
更具体地，对每个 \(x\in D(f(A))\)，有
\begin{equation}\label{b1:eq:29-spectral-graph-approximation}
 \|x-P_nx\|^2+
 \|f(A)x-f(A)P_nx\|^2
 =\int_{\{|f|>n\}}(1+|f|^2)\,\mathrm d\mu_x\longrightarrow0.
\end{equation}
特别地，\(\bigcup_{n\ge1}E([-n,n])H\) 在 \(D(A)\) 中关于图像范数稠密，且
\[
 \|x-E([-n,n])x\|^2+
 \|Ax-AE([-n,n])x\|^2
 =\int_{\{|t|>n\}}(1+t^2)\,\mathrm d\mu_x(t)\longrightarrow0
 \quad(x\in D(A)).
\]
\end{proposition}
\begin{proof}
上一小节已经证明 \(P_nH\subseteq D(f(A))\)。这些子空间递增，故其并为线性子空间。
由谱投影交换关系，
\(f(A)x-f(A)P_nx=f(A)(I-P_n)x\)。
\(I-P_n=E(\{|f|>n\})\)，所以分别对恒等乘数和 \(f\) 使用范数公式，
得到\eqref{b1:eq:29-spectral-graph-approximation}中的等式。
因为 \(1+|f|^2\) 对 \(\mu_x\) 可积，尾部积分趋于零。
平方和收敛等价于图像范数收敛，因而 \(\mathcal C_f\) 在定义域中关于图像范数稠密。
最后取 \(f(t)=t\)，得到关于 \(A\) 的结论。
\end{proof}

对一般 Borel 函数，采用 \(\{|f|\le n\}\) 作为截断集，保证每个投影的整个值域都在定义域中。
函数未必在实轴的有界区间上有界，所以不能无条件把它改为 \([-n,n]\)。
对于坐标函数，这两种截断恰好一致；所得谱子空间上的算子有界，
而图像范数逼近保证这些有界部分足以恢复原来的无界自伴算子。
